% RESULTADO_RECUPERADO: c40_termicas.tex del integral activo.
% Ampliación de prueba: conteo de modos, momentos térmicos, máximos
% espectrales y sustitución de la escala del calendario.
% Las hipótesis de realización se conservan explícitas.
\section{Rayonnement thermique et échelle d’information du calendrier}
\label{v:sec:radiacion-termica}

La statistique d’occupation fournit le lien entre les modes d’une
réalisation et leur énergie moyenne. Dans la réalisation électromagnétique,
la dispersion et la multiplicité de polarisation déterminent en outre la
densité spectrale. Leur composition permet d’obtenir les moments thermiques
et les maxima spectraux, en maintenant séparées trois opérations :
la production de l’échelle d’action, la représentation des modes et
l’évaluation de l’état thermique.

\subsection{Action, calendrier et température}

On conserve l’orientation \(a\) de l’horloge et on abrège
\(\hbar=\hbar_a\), \(c=c_{\rm int}=c_{\rm vac}\), avec
l’identification de \eqref{v:eq:identidad-velocidad} dans un même système de coordonnées.
L’énergie \(\varepsilon_{108}\) de cette section est
\(\epsilon_a\) de \eqref{v:ant:eq:ii-kb-aut-energia} ;
\(\Theta_{108}=\Theta_{\rm clk,a}\). Sa division par \(\log3\)
produit l’énergie par nat de \eqref{v:ant:eq:ii-kb-aut-energia-nat}.
La construction et la covariance de la paire thermique se trouvent dans
la section~\ref{v:ant:sec:ii-kb-aut-desarrollo}.

Nous écrivons \(h=2\pi\hbar>0\) pour la représentation d’action, \(t_0>0\)
pour l’unité temporelle de son système de coordonnées dimensionnelles et \(\ell_0=ct_0\) pour
la longueur correspondante. La cellule chronologique de longueur \(108\)
détermine l’énergie
\begin{equation}
 \varepsilon_{108}=\frac{h}{108t_0}
 =\frac{\hbar c}{L_\alpha}=Q_{L,a}=E_{P,a},\qquad
 k_B\Theta_{108}\log 3=Q_{L,a},\qquad
 \beta_{108}=\frac{108t_0\log3}{h}.
 \label{v:eq:energia-informacion-calendario}
\end{equation}
Ici \(\beta=(k_B\Theta)^{-1}\) est le paramètre de Gibbs. L’égalité
fixe le produit énergétique de la température et du transducteur
d’entropie. L’individualisation dimensionnelle de \(k_B\) et de
\(\Theta_{108}\) utilise la section thermique et ses bases. Le logarithme
naturel mesure l’information par branche tritique ; son facteur \(\log3\) transforme
la cardinalité de la cellule en entropie sans dimension.

Nous introduisons aussi les unités associées
\begin{equation}
 t_P=\frac{108}{2\pi_{\HMT}}t_0=\frac{L_\alpha}{c},\qquad
 \ell_P=ct_P=L_\alpha,\qquad
 E_P:=E_{P,a}=\frac{\hbar}{t_P}=Q_{L,a}=\varepsilon_{108}.
 \label{v:eq:unidades-termicas-calendario}
\end{equation}
Ces égalités précisent la signification des symboles dans cette section ;
leur évaluation dans un autre système de coordonnées conserve les identités entre quotients.
La longueur \(\ell_P=(54/\pi_{\HMT})\ell_0\) coïncide d’abord avec
la longueur structurelle \(L_\alpha\) de
\eqref{v:eq:longitud-alpha} ; l’identité gravitationnelle ultérieure
\(\ell_P^2=G_a\hbar_a/c^3\) est le corollaire
\eqref{v:eq:gravedad-unica-alpha}, non l’origine de cette échelle.

L’équivalence centrée
\eqref{v:eq:equivalencia-termica-centrada} et l’identité du transducteur
\eqref{v:eq:identidad-transductor-numeral} fixent, avant la réalisation
électromagnétique, l’échelle énergétique et sa lecture thermique. Le marqueur
$p_0$ de \eqref{v:eq:particion-portador-termico} conserve la normalisation
et l’origine lors du conditionnement d’une branche. Cette substitution détermine $\beta$ et
l’échelle d’énergie ; elle ne sélectionne pas rétrospectivement la dispersion ni
les deux polarisations spécifiées ci-après. C’est pourquoi les
lois de Planck, Stefan--Boltzmann et Wien restent des conséquences
ultérieures de la réalisation modale déclarée.

\subsection{Densité de modes et loi spectrale}

\begin{proposition}[Réalisation thermique des modes transversaux]
\label{v:prop:ley-espectral}
Considérons une réalisation libre dans une boîte périodique tridimensionnelle,
dont l’espace des excitations contient les vecteurs d’onde
\(\mathbf k\ne0\), avec deux modes transversaux par vecteur, une dispersion
\(\omega=c|\mathbf k|\), une énergie d’excitation \(\hbar\omega\) et
un potentiel chimique nul et \(\beta>0\). Le mode constant est maintenu
fixe comme donnée de fond. Dans la limite thermodynamique, la densité d’énergie
par fréquence angulaire est
\begin{equation}
 u_\omega(\Theta)=\frac{\hbar\omega^3}{\pi^2c^3}
 \frac{1}{e^{\beta\hbar\omega}-1},\qquad \omega>0.
 \label{v:eq:planck-angular}
\end{equation}
L’énergie est celle des excitations au-dessus du vide de cette
représentation.
\end{proposition}

\begin{proof}
Pour une boîte de volume \(V\), les vecteurs d’onde ont pour densité
\(V/(2\pi)^3\). Dans une couronne radiale d’épaisseur \(dk\), la densité
par volume, en incluant les deux polarisations, est
\[
 2\frac{4\pi k^2\,dk}{(2\pi)^3}=\frac{k^2}{\pi^2}\,dk
 =\frac{\omega^2}{\pi^2c^3}\,d\omega.
\]
Il s’agit de la densité limite des sommes sur le réseau périodique
sans le mode constant. Ce choix définit le produit de partition
avant la limite : chaque mode d’excitation a une énergie positive.
L’exclusion d’un seul point conserve la densité limite radiale.
L’occupation bosonique obtenue précédemment est
\(\bar n_\omega=(e^{\beta\hbar\omega}-1)^{-1}\).
Multiplier la densité par \(\hbar\omega\bar n_\omega\) démontre
\eqref{v:eq:planck-angular}. Près de zéro, l’intégrande énergétique est
d’ordre \(\omega^2\), et à l’infini, il décroît exponentiellement, de sorte
que son intégrale converge.
\end{proof}

La densité provient de la représentation tridimensionnelle spécifiée.
La multiplicité transversale utilisée ici a ce domaine physique ;
la dimension d’un espace harmonique sur le tore APP décrit,
séparément, la cohomologie de son complexe cellulaire.

% FORMALIZACION_NUEVA de la justificación del límite utilizado en
% 70_radiacion_termica.tex. No introduce un generador de constantes.
\subsection{Limite thermodynamique des sommes d’occupation}
\label{v:subsec:limite-termodinamico}

La densité radiale précédente peut être obtenue comme limite effective des
sommes de la réalisation périodique. L’estimation doit conserver le mode
constant fixé et contrôler uniformément aussi bien l’origine que la queue
spectrale. Les représentations d’action et de calendrier restent antérieures
à cette opération de réalisation.

\begin{lemma}[Convergence des trois observables thermiques]
\label{v:lem:limite-termodinamico}
Soient \(a,b>0\), \(L>0\), \(\Delta=2\pi/L\) et
\[
 F_N(r)=\frac{1}{e^{ar}-1},\qquad
 F_E(r)=\frac{br}{e^{ar}-1},\qquad
 F_Z(r)=-\log(1-e^{-ar})\quad(r>0).
\]
Pour chacune de ces fonctions, on a
\begin{equation}
 \lim_{L\to\infty}\frac{2}{L^3}
 \sum_{\nu\in\mathbb Z^3\setminus\{0\}}
 F\!\left(\frac{2\pi}{L}|\nu|\right)
 =\frac{1}{\pi^2}\int_0^\infty r^2F(r)\,dr.
 \label{v:eq:limite-modos-discretos}
\end{equation}
Les sommes et l’intégrale sont finies. La spécialisation
\(a=\beta\hbar c\), \(b=\hbar c\) fournit, respectivement,
les densités de nombre, d’énergie d’excitation et du logarithme de la fonction
de partition des modes transversaux.
\end{lemma}

\begin{proof}
Les trois fonctions sont continues et positives dans \((0,\infty)\).
Il existe des constantes \(C_0,C_1,d>0\), dépendant de \(a,b\) mais
indépendantes de \(\Delta\), telles que
\begin{equation}
 F(r)\leq\frac{C_0}{r}\quad(0<r\leq1),\qquad
 F(r)\leq C_1e^{-dr}\quad(r\geq1).
 \label{v:eq:envolventes-termicas}
\end{equation}
En effet, \(e^{ar}-1\geq ar\) contrôle \(F_N\) et \(F_E\)
près de l’origine ; en outre,
\(F_Z(r)=\log(1+F_N(r))\leq F_N(r)\leq1/(ar)\).
Pour \(r\geq1\), l’inégalité
\(-\log(1-y)\leq y/(1-y)\), avec \(y=e^{-ar}\), donne la borne
exponentielle de \(F_Z\). Les deux autres découlent de
\((e^{ar}-1)^{-1}\leq e^{-ar}/(1-e^{-a})\), en absorbant le facteur
\(r\) de \(F_E\) dans une exponentielle de taux inférieur. Ces bornes
prouvent déjà l’intégrabilité de \(r^2F(r)\).

Pour les sommes discrètes, nous regroupons les points entiers selon
\(m=\|\nu\|_\infty\geq1\). Chaque couche contient exactement
\[
 (2m+1)^3-(2m-1)^3=24m^2+2
\]
points et satisfait \(m\leq|\nu|\leq\sqrt3m\).
Fixons \(0<\delta\leq1\) et \(0<\Delta\leq1\).
Si \(\delta<\Delta\), il n’y a pas de points avec
\(0<\Delta|\nu|\leq\delta\). Sinon, en posant
\(N=\lfloor\delta/\Delta\rfloor\), la première borne de
\eqref{v:eq:envolventes-termicas} permet d’écrire
\begin{align*}
 \Delta^3\!\sum_{0<\Delta|\nu|\leq\delta}F(\Delta|\nu|)
 &\leq C_0\Delta^2\sum_{m=1}^{N}\frac{24m^2+2}{m}\\
 &\leq26C_0\Delta^2\sum_{m=1}^{N}m
 \leq26C_0\delta^2.
\end{align*}
Dans la deuxième ligne, on a utilisé \(m^{-1}\leq m\) puis
\(N(N+1)/2\leq N^2\) pour \(N\geq1\).
La contribution intégrale du même voisinage est aussi d’ordre
\(\delta^2\), car
\(\int_0^\delta r^2F(r)\,dr\leq C_0\delta^2/2\).

Pour \(R\geq1\), tout point avec \(\Delta|\nu|>R\) appartient
à une couche \(m>R/(\sqrt3\Delta)\). La seconde borne de
\eqref{v:eq:envolventes-termicas} donne
\begin{align*}
 \Delta^3\!\sum_{\Delta|\nu|>R}F(\Delta|\nu|)
 &\leq C_1e^{-dR/(2\sqrt3)}
 \Delta^3\sum_{m\geq1}(24m^2+2)e^{-d\Delta m/2}\\
 &\leq C_2e^{-dR/(2\sqrt3)},
\end{align*}
où \(C_2\) ne dépend pas de \(0<\Delta\leq1\). Pour vérifier
cette uniformité, on prend \(q=e^{-d\Delta/2}\) et on utilise
\[
 \sum_{m\geq1}q^m=\frac{q}{1-q},\qquad
 \sum_{m\geq1}m^2q^m=\frac{q(1+q)}{(1-q)^3}.
\]
La première identité provient de la série géométrique et la seconde de
l’application à deux reprises de \(q\,d/dq\), opération valable dans \(|q|<1\).
En outre,
\(1-e^{-d\Delta/2}\geq\Delta(1-e^{-d/2})\) par concavité.
Le facteur \(\Delta^3\) compense donc les dénominateurs des
deux expressions. La queue intégrale tend elle aussi vers zéro par la
borne exponentielle. Pour chaque \(\Delta>0\), ces mêmes estimations
prouvent la convergence de la somme infinie.

Dans la couronne compacte \(\delta\leq|k|\leq R\), la fonction
\(F(|k|)\) est continue et bornée. Les sommes de Riemann de maille
\(\Delta\) convergent vers son intégrale ; les deux sphères du bord
sont de mesure nulle. On fait d’abord \(\Delta\to0\), ensuite
\(\delta\to0\) et enfin \(R\to\infty\), en utilisant les bornes
uniformes précédentes. Il en résulte
\[
 \Delta^3\sum_{\nu\ne0}F(\Delta|\nu|)
 \longrightarrow \int_{\mathbb R^3}F(|k|)\,d^3k
 =4\pi\int_0^\infty r^2F(r)\,dr.
\]
Comme \(2/L^3=2\Delta^3/(2\pi)^3\), le coefficient radial final
est \(2\cdot4\pi/(2\pi)^3=1/\pi^2\), ce qui démontre
\eqref{v:eq:limite-modos-discretos}.
\end{proof}

Le contrôle de l’origine est effectué sur des excitations d’énergie positive.
Le mode constant reste fixé avant de former la fonction de partition :
son élimination ultérieure d’une intégrale ne remplacerait pas cette condition de
la représentation finie. La preuve établit le passage à la limite de la
réalisation déclarée et conserve la distinction entre cette réalisation et
la construction préalable de ses paramètres d’action et de calendrier.


\subsection{Moments thermiques et information}

\begin{lemma}[Moment bosonique]
\label{v:lem:momento-bosonico}
Pour tout réel \(s>1\),
\begin{equation}
 \int_0^\infty\frac{x^{s-1}}{e^x-1}\,dx
 =\Gamma(s)\zeta(s).
 \label{v:eq:momento-bosonico}
\end{equation}
\end{lemma}
\begin{proof}
Pour \(x>0\), \((e^x-1)^{-1}=\sum_{m\geq1}e^{-mx}\).
Tous les termes de la somme sont non négatifs ; Tonelli permet d’intervertir la somme
et l’intégrale. Le changement \(y=mx\) donne
\(\int_0^\infty x^{s-1}e^{-mx}\,dx=m^{-s}\Gamma(s)\).
La somme converge pour \(s>1\), et l’égalité annoncée en résulte.
\end{proof}

\begin{proposition}[Densités de nombre, d’énergie et d’entropie]
\label{v:prop:densidades-termicas}
Dans la réalisation de la proposition~\ref{v:prop:ley-espectral},
\begin{align}
 n_\gamma&=\frac{2\zeta(3)}{\pi^2c^3(\beta\hbar)^3},\\
 u_\gamma&=\frac{\pi^2}{15c^3\hbar^3\beta^4},\\
 s_\gamma&=\frac{4\pi^2 k_B}{45c^3(\beta\hbar)^3}.
 \label{v:eq:densidades-termicas}
\end{align}
À la température du calendrier, ces relations s’écrivent
\begin{equation}
 n_\gamma\ell_P^3=\frac{2\zeta(3)}{\pi^2(\log3)^3},\quad
 \frac{u_\gamma\ell_P^3}{E_P}=\frac{\pi^2}{15(\log3)^4},\quad
 \frac{s_\gamma\ell_P^3}{k_B}=\frac{4\pi^2}{45(\log3)^3}.
 \label{v:eq:densidades-calendario}
\end{equation}
\end{proposition}

\begin{proof}
L’intégrale de l’occupation multipliée par la densité de modes utilise
\eqref{v:eq:momento-bosonico} avec \(s=3\), et l’intégrale énergétique
l’utilise avec \(s=4\). On obtient \(\Gamma(3)=2\) et
\(\Gamma(4)\zeta(4)=6\pi^4/90=\pi^4/15\).
Pour l’entropie, on part du produit de partition des modes :
\[
 \frac{\log Z}{V}=-\frac{1}{\pi^2c^3}
 \int_0^\infty\omega^2
 \log(1-e^{-\beta\hbar\omega})\,d\omega
 =\frac{\pi^2}{45c^3(\beta\hbar)^3}.
\]
La dernière égalité s’obtient en intégrant par parties après le changement
\(x=\beta\hbar\omega\). Les termes de bord s’annulent :
\(x^3\log(1-e^{-x})\) tend vers zéro aux deux extrémités.
L’identité de Gibbs \(s_\gamma/k_B=\log Z/V+\beta u_\gamma\)
prouve la troisième expression. Enfin,
\(\beta_{108}\hbar=t_P\log3\) et \(\ell_P=ct_P\) donnent les trois
normalisations du calendrier.
\end{proof}

\subsection{Composition photonique du niveau \texorpdfstring{\(A_2\)}{A2} de Bendersky--Glaisher}
\label{v:subsec:bendersky-fotonica}

La densité de nombre de la proposition~\ref{v:prop:densidades-termicas}
admet une seconde lecture exacte qui conserve sa provenance spectrale.
Soit \(\mathcal Z_3\) la fonctionnelle scalaire normalisée produite par le
lecteur spectral tritique et reconnue ultérieurement comme la fonction zêta
dans la réalisation archimédienne. Soient
\[
 H_k=\sum_{j=1}^{k}\frac1j,\qquad H_0=0,
\]
et \(\mathsf B_j\) les nombres de Bernoulli avec la convention
\(\mathsf B_1=-1/2\). Pour éviter toute collision avec la notation des réseaux,
nous définissons les niveaux de Bendersky--Glaisher au moyen de
\begin{equation}
 A_k^{\mathrm{BG}}
 =\exp\!\left(
   \frac{H_k\mathsf B_{k+1}}{k+1}-\mathcal Z_3'(-k)
 \right),\qquad k\in\mathbb N_0.
 \label{v:eq:bendersky-bg}
\end{equation}
Le terme de Bernoulli fixe la normalisation du produit régularisé ;
il n’introduit pas de constante physique. La coordonnée
\(\pi_{\mathrm{HMT}}\) qui intervient dans l’équation fonctionnelle a déjà été
exprimée par l’état enrichi APP--TRIT--TPK avant cette lecture.

\begin{proposition}[Composition photonique du niveau deux]
\label{v:prop:bendersky-fotonica}
Dans la normalisation scalaire \(\mathcal Z_3=\zeta\),
\begin{equation}
 \log A_2^{\mathrm{BG}}
 =-\mathcal Z_3'(-2)
 =\frac{\mathcal Z_3(3)}{4\pi_{\mathrm{HMT}}^2}.
 \label{v:eq:bendersky-nivel-dos}
\end{equation}
Par conséquent, dans le système de coordonnées thermiques du calendrier,
\begin{align}
 n_\gamma\ell_P^3
  &=\frac{8\log A_2^{\mathrm{BG}}}{(\log3)^3}
   =-\frac{8\mathcal Z_3'(-2)}{(\log3)^3},
   \label{v:eq:numero-fotones-bendersky}\\
 \frac{u_\gamma}{n_\gamma k_B\Theta_{108}}
  &=\frac{\pi_{\mathrm{HMT}}^2}{120\log A_2^{\mathrm{BG}}},
   \label{v:eq:energia-media-bendersky}\\
 \frac{s_\gamma}{n_\gamma k_B}
  &=\frac{\pi_{\mathrm{HMT}}^2}{90\log A_2^{\mathrm{BG}}}.
   \label{v:eq:entropia-media-bendersky}
\end{align}
\end{proposition}

\begin{proof}
La forme complétée satisfait
\[
 \pi_{\mathrm{HMT}}^{-s/2}\Gamma(s/2)\mathcal Z_3(s)
 =\pi_{\mathrm{HMT}}^{-(1-s)/2}
  \Gamma((1-s)/2)\mathcal Z_3(1-s).
\]
En \(s=-2\), le zéro simple de \(\mathcal Z_3\) annule le pôle de
\(\Gamma(s/2)\). Comparer les termes constants, en utilisant
\(\Gamma(3/2)=\sqrt{\pi_{\mathrm{HMT}}}/2\) après la reconnaissance
archimédienne, donne
\[
 \mathcal Z_3'(-2)=-\frac{\mathcal Z_3(3)}{4\pi_{\mathrm{HMT}}^2}.
\]
Comme \(\mathsf B_3=0\), la définition~\eqref{v:eq:bendersky-bg}
produit la première égalité de~\eqref{v:eq:bendersky-nivel-dos}.
Substituer
\(\mathcal Z_3(3)=4\pi_{\mathrm{HMT}}^2\log A_2^{\mathrm{BG}}\)
dans~\eqref{v:eq:densidades-calendario} démontre
\eqref{v:eq:numero-fotones-bendersky}. Enfin,
\(k_B\Theta_{108}=E_P/\log3\) ; diviser les densités d’énergie et
d’entropie par la densité de nombre donne
\eqref{v:eq:energia-media-bendersky} et
\eqref{v:eq:entropia-media-bendersky}.
\end{proof}

L’évaluation \(A_2^{\mathrm{BG}}=1.0309167521973921\ldots\) sert de
test numérique des trois identités, non de donnée génératrice. Le développement
spectral de Bendersky--Glaisher contient la hiérarchie complète des dérivées
régularisées et leur correspondance avec les valeurs spectrales impaires ; la présente section
établit la composition réciproque qui utilise spécifiquement la densité
sans dimension de nombre photonique. On ne redéfinit pas \(k_B\), on n’identifie pas une densité à un
opérateur régularisé et on ne modifie aucune des densités avec
\(\zeta(3)\) déjà démontrées.


\subsection{Flux thermique et constante de Stefan--Boltzmann}

\begin{proposition}[Flux isotrope et normalisation d’action]
\label{v:prop:stefan}
Le flux hémisphérique de la réalisation isotrope est
\begin{equation}
 j_*(\Theta)=\sigma\Theta^4,\qquad
 \sigma=\frac{\pi^2 k_B^4}{60\hbar^3c^2}
 =\frac{2\pi^5k_B^4}{15h^3c^2}.
 \label{v:eq:stefan}
\end{equation}
À la température du calendrier,
\begin{equation}
 j_*(\Theta_{108})
 =\frac{2\pi^5}{15\,108^4(\log3)^4}
 \frac{h}{\ell_0^2t_0^2}.
 \label{v:eq:stefan-calendario}
\end{equation}
\end{proposition}
\begin{proof}
L’énergie isotrope par unité d’angle solide est
\(u_\gamma/(4\pi)\). Le flux normal à une surface intègre
\(c\cos\theta\) sur l’hémisphère :
\[
 \frac{cu_\gamma}{4\pi}
 \int_0^{2\pi}\!d\phi\int_0^{\pi/2}\cos\theta\sin\theta\,d\theta
 =\frac{cu_\gamma}{4}.
\]
La formule pour \(u_\gamma\), \(\beta^{-1}=k_B\Theta\) et
\(h=2\pi\hbar\) démontrent \eqref{v:eq:stefan}.
Substituer \eqref{v:eq:energia-informacion-calendario} et
\(c=\ell_0/t_0\) produit \eqref{v:eq:stefan-calendario}.
\end{proof}

La constante \(\sigma\) est, dans cette réalisation, le coefficient
qui compose l’occupation bosonique, la densité des modes transversaux et
la projection angulaire du flux. Sa puissance quatrième provient du moment
cubique de fréquence avec la différentielle d’intégration. La lecture
du calendrier exprime ce même flux au moyen de l’action, de la longueur,
du temps et de l’information tritique.

\newpage
\subsection{Maxima spectraux et loi de déplacement}

\begin{lemma}[Racine positive des maxima spectraux]
\label{v:lem:raiz-wien}
Pour tout entier \(n>1\), l’équation
\(n(1-e^{-x})=x\) a une unique solution \(x_n>0\).
La fonction \(f_n(x)=x^n/(e^x-1)\) a son unique maximum positif en
\(x_n\).
\end{lemma}
\begin{proof}
Soit \(g_n(x)=n(1-e^{-x})-x\). On a
\(g_n(0)=0\), \(g_n'(x)=ne^{-x}-1\) et
\(g_n''(x)=-ne^{-x}<0\). Le maximum de \(g_n\) est atteint en
\(\log n\), où sa valeur est \(n-1-\log n>0\), tandis que
\(g_n(n)=-ne^{-n}<0\). La monotonie stricte de chaque côté du maximum
prouve l’existence et l’unicité de la racine positive. Le signe de
\(f_n'(x)\) coïncide avec celui de \(g_n(x)\) ; en outre, \(f_n\) tend
vers zéro en zéro et à l’infini. L’unicité du maximum est ainsi prouvée.
\end{proof}

\begin{proposition}[Lectures en fréquence et en longueur d’onde]
\label{v:prop:wien}
Les maxima de luminance énergétique par fréquence et par longueur d’onde sont
\begin{equation}
 \nu_{\max}=\frac{x_3 k_B\Theta}{h},\qquad
 \lambda_{\max}=\frac{hc}{x_5 k_B\Theta}.
 \label{v:eq:wien}
\end{equation}
En particulier,
\begin{equation}
 \nu_{\max}(\Theta_{108})=
 \frac{x_3}{108t_0\log3},\qquad
 \lambda_{\max}(\Theta_{108})=
 \frac{108\log3}{x_5}\ell_0.
 \label{v:eq:wien-calendario}
\end{equation}
\end{proposition}
\begin{proof}
La conversion \(\omega=2\pi\nu\) et le facteur angulaire isotrope donnent
\[
 B_\nu=\frac{2h\nu^3}{c^2}\frac1{e^{\beta h\nu}-1}.
\]
Avec \(\nu=c/\lambda\), la densité transformée satisfait
\(B_\lambda=B_\nu|d\nu/d\lambda|\) ; par conséquent,
\[
 B_\lambda=\frac{2hc^2}{\lambda^5}
 \frac1{e^{\beta hc/\lambda}-1}.
\]
Les changements de variable respectifs conduisent à \(f_3\) et \(f_5\).
Le lemme précédent prouve \eqref{v:eq:wien} ; l’échelle du calendrier
donne \eqref{v:eq:wien-calendario}. Le jacobien de densité explique que
les deux maxima correspondent à des racines distinctes :
\(\nu_{\max}\lambda_{\max}=c x_3/x_5\).
\end{proof}

Le déplacement de Wien est défini structurellement par le point
stationnaire de la densité spectrale choisie. Le changement de variable
transforme aussi bien l’argument que la mesure. Cette précision situe
chaque constante spectrale dans son opération génératrice et conserve sa
signification physique lors du passage entre fréquence et longueur d’onde.
